\section{考察}
\label{sec:discussion}

\subsection{既存の枠組みとの関係}
\textbf{認知への力学系的アプローチ。}
概念を力学系のアトラクタとして扱うことには、コネクショニズムや身体性認知の研究に先例がある\citep[e.g.,][]{thelen1994dynamic}。
本枠組みは、VEDとIFGTの有効な動的慣例から構造的な動機を得て、明示的なトークン・場結合表式を提案する。
このsourceの再利用は、認知的アトラクタを上流の公理から導出したり、その機構を確立したりするものではない。

\textbf{分布意味論。}
分布意味論との比較は概念的なものである。
共起に関連するトークン空間の幾何を、CTが提案する履歴依存的なトークン・場幾何と対比する。
CTは\emph{使用}を重視するが、使用と\emph{共起}が経験的に排他的であるとも、この比較によって優れた意味表現が確立されるとも主張しない。
ここではコーパス、埋め込みの距離尺度、モデル間写像を与えない。

\textbf{予測処理と自由エネルギー。}
予測処理の記述は、認知を予測誤差の最小化として扱う\citep{friston2010free}。
ここでは概念的な比較を提案する。
$\sigma_{\mathrm{drive}}$は予測誤差の流入と並べて解釈でき、$\Phi$の井戸の深化は長期的な誤差最小化との類比で理解できる。
CTはこれらの量を予測誤差や自由エネルギーの変数と同一視しない。
ベイズ的生成モデルを採用するとは定めないが、その非採用は、より高い一般性や理論的包含の証明ではない。

\subsection{未解決問題}
\textbf{定量的対応。}
本枠組みは構造的記述を提案するが、畳み込み核$K$の具体的な形や、パラメータ$D$、$\mu$、$\alpha$、$\Gamma$、$\kappa_B$の値を指定しない。
これらを行動、言語、神経の経験的データから同定することは、主要な未解決問題である。

\textbf{トポロジーと概念の個体化。}
本枠組みは概念を$\Phi$の局所的安定領域として定義する。
しかし、ある領域が関連する概念の集まりではなく単一の概念と数えられる条件は、まだ指定されていない。
概念の個体化にはトポロジー的基準が必要である。

\textbf{マルチエージェントの同期ダイナミクス。}
社会的地形は定性的に記述している。
共有地形が安定に形成される条件や破綻する条件を含め、マルチエージェントのトークン・場相互作用を形式的に扱う方法は今後の課題である。

\textbf{神経実装との接続。}
場の量$I$、$\Phi$、$\JI$と、測定可能な神経量（発火率、結合重み、振動ダイナミクス）との対応は、経験的神経科学との協働を必要とする主要な未解決問題である。
